\section{开场：这门课在讲什么}

CS336 的第一讲不是直接进入模型细节，而是先回答一个更根本的问题：为什么要在 2025 年再开一门“从零构建语言模型”的课程。
讲师一开始就把课程定位得很清楚：这门课不是教你追逐最新的热点技巧，也不是让你依赖一个黑箱 API 直接调用，而是要把语言模型整个栈拆开，重新理解其数据、系统、架构和训练逻辑。

\begin{importantbox}{本讲的核心信息}
\begin{itemize}
    \item 课程目标不是“会用模型”，而是“理解模型如何被构建出来”。
    \item 课程强调完整栈：数据、tokenization、模型结构、训练、系统、规模律、对齐。
    \item 课程的姿态不是反对抽象，而是警惕抽象泄漏，并在必要时回到机制层面。
\end{itemize}
\end{importantbox}

讲师给出的课程口号可以概括成一句话：
\begin{quote}
理解它，必须先构建它。
\end{quote}

